\chapter{Jacques Tits (2008)}

\section{Biographical Sketch}

Jacques Tits was born on 12 August 1930 in Uccle, Belgium. He studied at the Universit\'e Libre de Bruxelles, where he earned his PhD in 1954 under the supervision of Georges Lema\^itre (the physicist who proposed the Big Bang theory). He held positions at the Universit\'e Libre de Bruxelles, the University of Bonn, and the Coll\`ege de France. Tits received the Abel Prize (jointly with Thompson) in 2008, the Wolf Prize in 1993, and the Cantor Medal in 1996.

\section{High-Level View for a General Audience}

\subsection{What Did Tits Do?}

Imagine you have a collection of shapes---triangles, squares, pentagons---and you want to understand all the ways you can rotate and reflect them so they land back on themselves. For a square, there are 8 such motions: 4 rotations and 4 reflections. For a regular pentagon, there are 10. This collection of motions is called a \emph{symmetry group}.

Now imagine the ultimate challenge: classify all possible symmetry groups, in any number of dimensions, with any number of elements. This is the problem of classifying \emph{finite simple groups}---the atomic building blocks of all finite symmetry groups. Tits played a central role in this monumental achievement.

\paragraph{Tits' Contribution:} Tits created the theory of \emph{Buildings\index{Tits buildings}}, a geometric structure that encodes the symmetries of exceptional Lie groups\index{Lie groups}. This theory allowed mathematicians to understand the deep geometry underlying groups that had previously seemed abstract and mysterious. Tits also discovered a new sporadic simple group (the Tits group), formulated the Tits alternative in geometric group theory, and provided a uniform classification of groups of Lie type through his theory of BN-pairs.

\subsection{Why Does It Matter?}

Symmetry groups are everywhere in mathematics and physics:

\begin{itemize}
\item They describe the structure of molecules in chemistry.
\item They are the foundation of particle physics (the Standard Model).
\item They are used in cryptography to secure communications.
\item They describe the symmetries of crystal structures.
\item They appear in art, architecture, and music.
\end{itemize}

The classification of finite simple groups is one of the greatest intellectual achievements of all time, comparable to the classification of chemical elements in the periodic table.

\section{The Origin of the Problem}

\subsection{Group Theory and Symmetry}

Group theory was born in the early nineteenth century through the work of Galois, Abel, Cauchy, and others. A \emph{group} is a set equipped with an associative binary operation, an identity element, and inverses. Groups capture the abstract notion of symmetry.

\begin{knowledgebridge}{What Is a Group?}
A \emph{group} is a set of actions that can be combined and reversed. For a square, the set of rotations and reflections forms a group: you can rotate by $90^\circ$ (an action), do it twice (combine actions), or undo it (take the inverse). Formally, a group requires three ingredients: an \emph{associative} operation (combining actions), an \emph{identity} element (doing nothing), and \emph{inverses} (undoing actions). Groups appear everywhere: the symmetries of a molecule, the set of moves on a Rubik's cube, and the fundamental laws of particle physics.
\end{knowledgebridge}

A group is \emph{simple} if it has no nontrivial normal subgroups---it cannot be decomposed into smaller pieces. Simple groups are the atoms of group theory: every finite group can be built from simple groups by repeated extensions (the Jordan--H\"older theorem).

\begin{knowledgebridge}{What Makes a Group Simple?}
A group is \emph{simple} if it cannot be broken into smaller, nontrivial pieces---like a chemical element that cannot be split by ordinary means. Just as every molecule is built from elements in the periodic table, every finite group can be built from simple groups (the Jordan--H\"older theorem). The classification of finite simple groups is the mathematical equivalent of listing all chemical elements.
\end{knowledgebridge}

\subsection{The Search for Simple Groups}

By the 1950s, mathematicians had identified several infinite families of finite simple groups:

\begin{itemize}
\item \textbf{Cyclic groups} $\mathbb{Z}_p$ of prime order.
\item \textbf{Alternating groups} $A_n$ for $n \geq 5$.
\item \textbf{Groups of Lie type}: families of matrix groups over finite fields, including $PSL(n, q)$, $PSp(2n, q)$, and the exceptional groups $G_2(q)$, $F_4(q)$, $E_6(q)$, $E_7(q)$, $E_8(q)$.
\end{itemize}

In addition, a handful of "sporadic" simple groups had been discovered, including the Mathieu groups $M_{11}$, $M_{12}$, $M_{22}$, $M_{23}$, $M_{24}$.

\subsection{The Classification Problem}

The central problem was: are there any more finite simple groups beyond these families? And can we prove that the list is complete?

This became known as the \emph{classification of finite simple groups} (CFSG), one of the most ambitious projects in the history of mathematics. The complete classification was announced in 1983 (with final corrections in 2004). It states that every finite simple group is:

\begin{enumerate}
\item A cyclic group of prime order;
\item An alternating group $A_n$ ($n \geq 5$);
\item A group of Lie type (including the Tits group);
\item One of 26 sporadic groups.
\end{enumerate}

\begin{analogy}{The Periodic Table of Symmetry}
Classifying finite simple groups is like creating a periodic table for symmetry. Just as Mendeleev discovered that chemical elements fall into families with predictable properties, mathematicians found that all finite simple groups belong to a few infinite families (cyclic, alternating, Lie type) plus 26 exceptional ``sporadic'' groups. The sporadic groups are like the noble gases---they do not quite fit the patterns of the families, yet they are essential for a complete picture.
\end{analogy}

\begin{figure}[htbp]
\centering
\begin{tikzpicture}[
    node distance=4mm and 4mm,
    box/.style={draw, thick, rounded corners, align=center, minimum width=2.8cm, minimum height=0.7cm, font=\small},
    titlebox/.style={draw, thick, rounded corners, align=center, font=\bfseries\small, minimum width=2.8cm},
    legend/.style={font=\footnotesize, align=center}
]

% Title
\node[font=\Large\bfseries, abelblue] at (6.5, 9.2) {Classification of Finite Simple Groups};

% --- Section 1: Cyclic ---
\node[box, fill=abelblue!15, draw=abelblue] (cyclic) at (2, 7.5) {$\mathbb{Z}_p$\\ {\footnotesize (prime order)}};
\node[legend, above=0pt of cyclic] {\textbf{Cyclic Groups}};

% --- Section 2: Alternating ---
\node[box, fill=abelred!15, draw=abelred, minimum width=3.5cm] (alt) at (6.5, 7.5) {$A_n$ ($n \ge 5$)\\ {\footnotesize alternating groups}};
\node[legend, above=0pt of alt] {\textbf{Alternating Groups}};

% --- Section 3: Lie Type ---
\node[titlebox, fill=abelgreen!15, draw=abelgreen, minimum width=8cm, minimum height=0.6cm] (lie) at (6.5, 5.8) {Groups of Lie Type};
\node[box, fill=abelgreen!8, draw=abelgreen, minimum width=7.5cm, minimum height=1.2cm] (classical) at (4.2, 4.5) {{\bf Classical:}\\ $PSL(n,q)$, $PSp(2n,q)$, $PSU(n,q)$, $PO(n,q)$};
\node[box, fill=abelgreen!8, draw=abelgreen, minimum width=7.5cm, minimum height=1.2cm] (exceptional) at (8.8, 4.5) {{\bf Exceptional:}\\ $G_2$, $F_4$, $E_6$, $E_7$, $E_8$};
\draw[abelgreen, thick] (lie.south) -- (classical.north);
\draw[abelgreen, thick] (lie.south) -- (exceptional.north);

% --- Section 4: Sporadic ---
\node[titlebox, fill=abelgray!15, draw=abelgray, minimum width=5cm] (spor) at (6.5, 2.5) {26 Sporadic Groups};
\node[box, fill=abelred!20, draw=abelred, minimum width=4cm, minimum height=0.8cm] (monster) at (6.5, 1.2) {{\LARGE\bfseries Monster}\\ $|M| \approx 8\times 10^{53}$};
\draw[abelgray, thick] (spor.south) -- (monster.north);

% Note about Tits group
\node[box, fill=yellow!15, draw=abelgray, minimum width=5cm] (tits) at (6.5, -0.3) {${}^2F_4(2)'$ (Tits group) is sporadic};
\draw[dashed, abelgray, thick] (monster.south) -- (tits.north);

% Connecting brackets on left
\draw[thick, abelblue] (0.2, 7.85) -- (0.2, 7.15);
\draw[thick, abelred] (0.2, 7.15) -- (0.2, 6.85);
\draw[thick, abelgreen] (0.2, 6.85) -- (0.2, 3.5);
\draw[thick, abelgray] (0.2, 3.5) -- (0.2, 0.4);

\end{tikzpicture}
\caption{The classification of finite simple groups: every finite simple group belongs to one of four families — cyclic groups of prime order, alternating groups of degree at least 5, groups of Lie type (classical or exceptional), or one of the 26 sporadic groups. The Tits group ${}^2F_4(2)'$, though closely related to Lie type, is classified as sporadic.}
\label{fig:finite_simple_groups_tits}
\end{figure}

\section{Deep Insight into the Problem}

\subsection{Tits' Buildings}\index{Tits buildings}

Jacques Tits introduced the concept of a \emph{building} in the 1960s. A building is a simplicial complex that encodes the structure of a group of Lie type. It consists of:
\begin{itemize}
\item \textbf{Apartments}: Subcomplexes isomorphic to the Coxeter complex of a Weyl group.
\item \textbf{Chambers}: Maximal simplices in the building.
\item \textbf{Gallery Connections}: A system of adjacency between chambers.
\end{itemize}

Buildings provide a geometric realization of the BN-pair structure of a group of Lie type. The building of $SL(n, q)$, for example, is the flag complex of subspaces of $\mathbb{F}_q^n$.

\begin{primer}{Tits Buildings}
A Tits \emph{building} is a geometric structure that encodes the symmetries of a group. Imagine building a crystal lattice: the arrangement of atoms reflects the underlying symmetry group. A building generalizes this idea: it is a collection of simplices (points, lines, triangles) arranged into \emph{apartments} (flat regions) that overlap in a controlled way. For the group $SL(n,q)$, the building is the set of all subspaces of $\mathbb{F}_q^n$, where the chambers are maximal chains of subspaces. Buildings give a geometric language for understanding groups that otherwise seem purely algebraic.
\end{primer}

\subsection{Tits' Geometries and the Tits Group}

Tits discovered a new sporadic simple group during his work on geometries over finite fields. The \emph{Tits group} $^2F_4(2)'$ is a simple group of order 17,971,200 that arises from the geometry of the Ree groups $^2F_4(q)$.

More generally, Tits showed that many of the sporadic groups discovered by Fischer, Janko, Conway, and others are related to geometries that can be understood using his building theory.

\subsection{BN-Pairs (Tits Systems)}

Tits introduced the notion of a \emph{BN-pair} (also called a Tits system) as an axiomatic framework for groups of Lie type. A BN-pair consists of two subgroups $B$ and $N$ of a group $G$ satisfying certain axioms that capture the essential structure of a group with a Bruhat decomposition. This formalism provides a uniform way to study all groups of Lie type, both finite and infinite, and leads directly to the construction of the associated building.

\section{Biggest Contributions and New Tools}

\subsection{The Theory of Buildings}

Tits' theory of buildings is a monumental achievement:

\begin{itemize}
\item \textbf{Spherical Buildings}: Associated to isotropic algebraic groups over fields, these are the most important class. The classification of spherical buildings of rank at least 3 is a landmark result.

\item \textbf{Affine Buildings}: Associated to algebraic groups over local fields, these buildings are fundamental in the study of $p$-adic groups and the Langlands program.

\item \textbf{Twin Buildings}: Tits introduced twin buildings as a geometric framework for Kac--Moody groups.
\end{itemize}

\subsection{The Tits Alternative}

A fundamental result in geometric group theory, the Tits alternative states that every finitely generated linear group either contains a non-abelian free subgroup or is virtually solvable. This dichotomy has profound consequences for the structure of linear groups and their geometry.

\subsection{The Classification of Groups of Lie Type}

Tits' work on groups generated by root subgroups (the "BN-pair" axioms) provided a uniform approach to all groups of Lie type. His classification of irreducible spherical buildings of rank at least 3 gave an alternative proof of the classification of simple algebraic groups over arbitrary fields.

\section{Impact of the Discovery in Science and Technology}

\subsection{Impact on Mathematics}

\begin{whyitmatters}
The classification of finite simple groups---to which Tits contributed decisively---is one of the greatest intellectual achievements in mathematics, comparable to the classification of chemical elements or the mapping of the human genome. Tits' theory of buildings provided the geometric language for understanding the groups of Lie type, and his BN-pair axioms unified the entire theory of algebraic groups over arbitrary fields.
\end{whyitmatters}

Tits' work transformed group theory and its connections:

\begin{itemize}
\item \textbf{Geometric Group Theory}: Tits' work on buildings and the Tits alternative laid foundations for geometric group theory, which studies groups through their actions on metric spaces.

\item \textbf{Representation Theory}: Tits' geometric approach to representations has shaped modern representation theory.

\item \textbf{Algebraic Geometry}: The theory of buildings is essential for understanding algebraic groups over arbitrary fields and their geometry.

\item \textbf{Lie Theory}: Tits' classification of spherical buildings provided a geometric proof of the classification of simple algebraic groups.
\end{itemize}

\subsection{Impact on Physics}

\begin{itemize}
\item \textbf{Particle Physics}: The classification of simple groups explains the symmetry groups of fundamental interactions. The Standard Model's gauge group $SU(3) \times SU(2) \times U(1)$ is a product of simple groups and abelian factors.

\item \textbf{String Theory}: Exceptional groups, including $E_8$ which is related to Tits' building theory, appear in string theory and M-theory.

\item \textbf{Condensed Matter Physics}: Symmetry groups are used to classify phases of matter, and the classification of simple groups provides the complete list of possible symmetry types.
\end{itemize}

\subsection{Impact on Technology}

\begin{itemize}
\item \textbf{Cryptography}: Group theory is used in cryptography, including the use of group structures in elliptic curve\index{Elliptic curves} cryptography and lattice-based cryptography.

\item \textbf{Coding Theory}: Group-theoretic codes, including group codes and Reed--Muller codes, use the classification of simple groups.

\item \textbf{Quantum Computing}: The theory of quantum error-correcting codes uses group theory, and the classification of simple groups provides constraints on possible code structures.
\end{itemize}

\section{Current Developments}

\subsection{Buildings and the Langlands Program}

Tits' buildings play a central role in the Langlands program:

\begin{itemize}
\item \textbf{$p$-adic Groups}: The theory of buildings is essential for studying representations of $p$-adic groups and the local Langlands correspondence.

\item \textbf{Rapoport--Zink Spaces}: These spaces, which parametrize $p$-adic period domains, are studied using building theory.

\item \textbf{Geometric Satake}: The geometric Satake correspondence, which is central to the geometric Langlands program, involves the affine Grassmannian, which is related to the building of the loop group.
\end{itemize}

\subsection{The Second-Generation Classification}

Work continues on simplifying and verifying the classification of finite simple groups:

\begin{itemize}
\item \textbf{The GLS Project}: Gorenstein, Lyons, and Solomon are writing a series of volumes providing a complete and unified proof of the classification.

\item \textbf{Aschbacher's Contributions}: Aschbacher has made major contributions to simplifying portions of the classification proof.

\item \textbf{Computer Verification}: Formal verification of parts of the classification using proof assistants is an ongoing project.
\end{itemize}

\section{Conclusion}

Jacques Tits transformed our understanding of symmetry. His theory of buildings provided the geometric language for understanding the groups of Lie type, his BN-pair axioms unified the theory of algebraic groups, and his Tits alternative became a cornerstone of geometric group theory.

The Abel Prize citation recognized Tits (jointly with Thompson) "for their profound achievements in algebra and in particular for shaping modern group theory." This captures only a fraction of their impact: Tits' work has shaped not just algebra, but geometry, representation theory, number theory, and mathematical physics.

\section{Behind the Breakthrough: The Human Story}
Tits' collaboration with Armand Borel produced the foundational Borel--Tits paper on reductive groups, and Tits was known for his remarkable geometric intuition. Thompson's Feit-Thompson proof (establishing the Odd Order Theorem) was 255 pages long, occupying an entire issue of the Pacific Journal of Mathematics---at the time completely unprecedented, it paved the way for the classification of finite simple groups alongside Tits' geometric contributions.

\section{Pedagogical Metaphors and Lineage}
\subsection{Signature Metaphor}
``Constructing the geometric 'skeletons' (buildings) that support the algebraic symmetries of groups of Lie type.''

\subsection{Mathematical Lineage}
Tits was advised by Georges Lema\^itre (the physicist who proposed the Big Bang theory). Tits in turn advised numerous mathematicians who continued the development of building theory and algebraic groups.

\section{Applied Science and Computational Algorithms}
\subsection{Key Takeaways for Applied Science}
Group theory and buildings are used in crystallography (analyzing atomic configurations), structural chemistry (molecular symmetries), and in coding theory for designing algebraic error-correcting codes.

\subsection{Algorithms and Numerical Implementations}
Coset enumeration algorithms (such as the Todd-Coxeter algorithm) and the Schreier-Sims algorithm for permutation groups (implemented in computer algebra systems like GAP and Magma) are used to compute simple group representations.

\section{The Research Frontier}
Understanding the geometry of buildings associated to infinite-dimensional Kac-Moody groups, and completing the second-generation proof of the Classification of Finite Simple Groups.

\section{Guided Exercises and Challenges}
\begin{enumerate}[label=\textbf{\arabic*.}]
  \item Define a Tits building and describe its structure in terms of apartments and chambers.
  \item What is a BN-pair (or Tits system) and how does it give rise to a building?
  \item State the Tits alternative and explain its significance in geometric group theory.
\end{enumerate}

\section{Curriculum Map and Further Reading}
For students wishing to delve deeper into the mathematical machinery underlying these breakthroughs, the following learning path is recommended:
\begin{itemize}
  \item Jacques Tits, \emph{Buildings of Spherical Type and Finite BN-Pairs}, Lecture Notes in Mathematics, Springer.
  \item Michael Aschbacher, \emph{The Finite Simple Groups}, Cambridge University Press.
\end{itemize}

\section*{Bibliography}

\begin{enumerate}
\item J. Tits, "Buildings of Spherical Type and Finite BN-Pairs," Springer, 1974.
\item J. Tits, "Th\'eorie des groupes," Annuaire du Coll\`ege de France 1980--1990.
\item J. Tits, "Le probl\`eme des mots dans les groupes de Coxeter," Symposia Mathematica 1 (1969), 175--185.
\item A. Borel and J. Tits, "Groupes r\'eductifs," Publications Math\'ematiques de l'IH\'ES 27 (1965), 55--151.
\item J. Tits, "Groupes de Lie et leurs g\'eom\'etries," Annales de l'Institut Fourier 4 (1952), 1--44.
\item J. Tits, "G\'eom\'etries poly\'edriques et groupes simples," Atti del convegno di geometria combinatoria e sue applicazioni (1971), 371--387.
\item W. Feit and J. G. Thompson, "Solvability of groups of odd order," Pacific Journal of Mathematics 13 (1963), 775--1029.
\item D. Gorenstein, "Finite Simple Groups: An Introduction to Their Classification," Plenum, 1982.
\item M. Aschbacher, "The Classification of Finite Simple Groups," American Mathematical Society, 2004.
\end{enumerate}
